\documentclass[11pt]{article}

\usepackage[utf8]{inputenc}
\usepackage[T1]{fontenc}
\usepackage{lmodern}
\usepackage[margin=1in]{geometry}
\usepackage{booktabs}
\usepackage{amsmath}
\usepackage{microtype}

\title{Natural Language Processing - Ex4}
\author{Nurit Tolkowsky and Shaul Tolkowsky}
\date{}

\begin{document}
\maketitle

\section{MST Parser}

\subsection*{1.4 \quad Evaluation (UAS)}

We trained the averaged perceptron on the first 90\% of the
\texttt{dependency\_treebank} parsed sentences (3522 sentences), for 2
iterations with learning rate 1, traversing the training instances in a
random order. Edge scores are the dot product of the weight vector with the
word-bigram and POS-bigram feature function; inference uses the
Chu-Liu/Edmonds algorithm (\texttt{cle\_min}) applied to negated scores
(minimum arborescence $=$ maximum spanning tree). Evaluated on the last 10\%
of the corpus (392 sentences), the mean unlabeled attachment score (UAS) of
the learned averaged weight vector is:
\[
  \text{UAS (averaged perceptron)} = 0.1768 .
\]

\section{Attention-Based Parsing}

\subsection*{2.2(b) \quad UAS per BERT layer (\texttt{head\_mode = "mean"})}

Using \texttt{get\_attention\_matrix}, \texttt{attn\_to\_arc\_scores} (which
negates the attention weights for \texttt{cle\_min} and never creates an edge
into ROOT, since ROOT can never be a dependent), and \texttt{cle\_min}, we
parsed every sentence in the test set. The mean UAS for each setting is:

\begin{center}
\begin{tabular}{lc}
\toprule
Setting (\texttt{head\_mode = mean}) & UAS \\
\midrule
\texttt{layer = 0}\ \ (first BERT layer)  & 0.230 \\
\texttt{layer = 5}\ \ (middle BERT layer) & 0.138 \\
\texttt{layer = 11} (last BERT layer)     & 0.120 \\
\bottomrule
\end{tabular}
\end{center}

\paragraph{Discussion - which setting works best and why.}
The first layer (\texttt{layer = 0}) gives the highest UAS. Because
\texttt{head\_mode="mean"} averages all heads, the score reflects each
layer's average attention pattern: in the lowest layer this pattern is
dominated by local, positional behaviour (words attend mostly to their
immediate neighbours). Since a large fraction of dependency arcs connect
nearby words, this surface locality lines up with the gold tree better than
the more diffuse, semantically-oriented attention of the middle and upper
layers, where the few genuinely syntactic heads are washed out by averaging.

This may seem to contradict the intuition stated in the exercise - that,
since lower layers are closer to surface syntax and upper layers to
semantics, the middle layers should be most useful for parsing. A likely
reason it does not hold here is that we average all heads together
(\texttt{head\_mode="mean"}): if only a few individual heads track syntax,
averaging dilutes them, so the score is dominated by each layer's overall
attention pattern. In the lowest layer that pattern is strongly local, and
since many dependencies are between nearby words, \texttt{layer = 0} ends up
winning.

\subsection*{2.3 \quad Comparison with the MST parser}

\begin{center}
\begin{tabular}{lc}
\toprule
Parser & UAS \\
\midrule
Perceptron MST (averaged weights)        & 0.1768 \\
BERT attention, best setting (\texttt{layer = 0}) & 0.230 \\
\bottomrule
\end{tabular}
\end{center}

\paragraph{Data requirements.}
The perceptron needs labeled, dependency-annotated training data plus a
hand-designed feature function, whereas the attention-based parser needs no
task-specific training at all - it reuses attention weights from a model
pre-trained only on masked language modeling.

\paragraph{Inference cost.}
The perceptron is far cheaper, scoring each candidate edge with a sparse dot
product, while the attention parser requires a full BERT forward pass per
sentence.

\paragraph{Accuracy.}
Here the best attention setting (0.230) slightly outperforms our 2-iteration
perceptron (0.1768), which is notable given that BERT never saw a parse tree.
Part of the perceptron's low score is that its feature function is
deliberately simple: every feature looks only at a pair - the two endpoints
of a single edge (one word bigram and one POS bigram) - with no higher-order
features (such as word triplets, head-dependent distance, or combined
word+POS features) and no surrounding context. The word-bigram features are
also extremely sparse - most word pairs seen at test time never appeared in
training - so they barely generalize, and with only 2 training iterations the
model stays underfit. BERT, by contrast, encodes each word using the whole
sentence, so even its raw attention carries far richer structural
information.

\paragraph{Why untrained attention correlates with syntax.}
Accurately predicting a masked word requires implicitly modeling grammatical
relationships - which words govern or modify which - so some attention heads
end up tracking head-dependent structure as a by-product of language
modeling.

\end{document}
